% Metódy inžinierskej práce -- predbežná verzia projektu
% Formulár VV-A až VV-E žiadosti pre aplikovaný výskum (podľa vzoru APVV VV 2024, Príloha 5)
\documentclass[10pt,a4paper]{article}

\usepackage[slovak]{babel}
\usepackage[T1]{fontenc}
\usepackage[utf8]{inputenc}
\usepackage{lmodern}
\usepackage[margin=2cm]{geometry}
\usepackage{array}
\usepackage{tabularx}
\usepackage{longtable}
\usepackage{booktabs}
\usepackage{xcolor}
\usepackage{fancyhdr}
\usepackage[hidelinks]{hyperref}

\definecolor{formbg}{gray}{0.92}
\newcommand{\akronym}{AI-PENTEST}
\newcommand{\doplnit}[1]{\textcolor{gray}{[#1]}}
\newcommand{\polozka}[3]{\textbf{#1} & \cellcolor{formbg}\small #2 & #3 \\ \hline}
\usepackage{colortbl}

\pagestyle{fancy}
\fancyhf{}
\lhead{\small VV 2024 \quad Aplikovaný výskum / Applied research}
\rhead{\small Akronym: \akronym}
\cfoot{\small Strana/Page: \thepage}
\setlength{\headheight}{14pt}

\newenvironment{formular}[2]{%
  \section*{#1 \quad {\normalsize #2}}
  \noindent\renewcommand{\arraystretch}{1.35}%
  \begin{longtable}{|p{0.6cm}|p{4.6cm}|p{10.3cm}|}\hline
}{%
  \end{longtable}
}

\begin{document}

\begin{center}
  {\Large\bfseries Žiadosť o poskytnutie prostriedkov na riešenie projektu}\\[2pt]
  {\large Aplikovaný výskum -- formulár VV-A až VV-E}\\[4pt]
  {\small Predbežná verzia, Metódy inžinierskej práce, ak. rok 2026/27}
\end{center}

% ---------------------------------------------------------------
\begin{formular}{VV-A1}{Základné informácie o projekte}
\polozka{01}{Evidenčné číslo projektu}{APVV-24-0000}
\polozka{02}{Dátum podania}{1.~11.~2026}
\polozka{03}{Názov projektu}{Inteligentný systém na automatizované testovanie bezpečnosti webových aplikácií pomocou veľkých jazykových modelov}
\polozka{03}{Názov projektu v~angličtine}{An Intelligent System for Automated Security Testing of Web Applications Using Large Language Models}
\polozka{04}{Akronym projektu}{\akronym}
\polozka{05}{Odbor vedy a techniky}{Technické vedy -- informatika}
\polozka{06}{Charakter výskumu}{Aplikovaný výskum}
\polozka{07}{Začiatok riešenia projektu}{1.~7.~2027}
\polozka{08}{Koniec riešenia projektu}{30.~6.~2030}
\polozka{09}{Anotácia (max. 2\,000 znakov)}{%
Súčasné veľké jazykové modely dokážu vykonávať zložité úlohy súvisiace s~analýzou webových aplikácií, no chýba štandardizované prostredie na bezpečné a objektívne porovnávanie ich schopností v~oblasti penetračného testovania. Projekt \akronym{} preto navrhuje vytvoriť Web Pentest Harness -- kontrolované prostredie, v~ktorom AI agent dostane webovú aplikáciu a pokúsi sa samostatne vykonať jej bezpečnostné testovanie. Harness poskytuje agentovi izolované testovacie prostredie a nástroje na interakciu s~aplikáciou, zaznamenáva všetky jeho akcie, overuje existenciu nájdenej zraniteľnosti kontrolovaným proof-of-concept postupom, vyhodnocuje úspešnosť agenta a zbiera výsledky na porovnanie viacerých AI modelov. Jadrom riešenia je kontajnerizovaná sada zámerne zraniteľných webových aplikácií so známymi triedami chýb (SQL injection, XSS, obchádzanie autentifikácie a autorizácie, SSRF, IDOR, nesprávna konfigurácia), doplnená anotovanou databázou referenčných zraniteľností a hodnotiacim modulom, ktorý počíta presnosť, úplnosť, mieru falošne pozitívnych nálezov a čas do nájdenia a overenia zraniteľnosti. Výstupom bude otvorená infraštruktúra (harness a benchmark úloh rôznej obtiažnosti) a empirické porovnanie architektúr AI agentov z~hľadiska efektivity, spoľahlivosti a typických zlyhaní.}
\polozka{10}{Žiadateľská organizácia}{Slovenská technická univerzita v~Bratislave, Fakulta informatiky a informačných technológií}
\polozka{11}{Požadované finančné prostriedky z~APVV (v~EUR)}{240\,000}
\polozka{12}{Spolufinancovanie projektu (v~EUR)}{0}
\polozka{13}{Celkové náklady na projekt (v~EUR)}{240\,000}
\end{formular}

% ---------------------------------------------------------------
\begin{formular}{VV-A2}{Žiadateľ -- základné informácie o riešiteľskej organizácii}
\polozka{01}{Názov organizácie}{Slovenská technická univerzita v~Bratislave}
\polozka{02}{Skrátený názov}{STU}
\polozka{03}{Adresa organizácie}{Vazovova 5, 812 43 Bratislava; \newline pracovisko: FIIT STU, Ilkovičova 2, 842 16 Bratislava}
\polozka{04}{IČO}{00397687}
\polozka{05}{Príslušnosť k~rezortu}{Ministerstvo školstva, výskumu, vývoja a mládeže SR}
\polozka{06}{Forma hospodárenia}{Verejná vysoká škola}
\polozka{07}{Kontaktná osoba / e-mail}{Andrii Tsybulskyi, \texttt{xtribulskyi@stuba.sk}}
\end{formular}

\noindent\emph{Spoluriešiteľská organizácia:} projekt nemá spoluriešiteľskú organizáciu.

% ---------------------------------------------------------------
\section*{VV-A3 \quad {\normalsize Zoznam riešiteľov}}
\renewcommand{\arraystretch}{1.3}
\begin{tabularx}{\textwidth}{|>{\raggedright\arraybackslash}X|l|l|r|}
\hline
\rowcolor{formbg}\textbf{Meno a priezvisko} & \textbf{Pracovné zaradenie} & \textbf{IČO org.} & \textbf{Počet hodín} \\ \hline
Andrii Tsybulskyi (zodpovedný riešiteľ) & výskumný pracovník & 00397687 & 2\,400 \\ \hline
Kyrylo Bukhanets & výskumný pracovník & 00397687 & 2\,000 \\ \hline
Artur Toba & výskumný pracovník & 00397687 & 2\,000 \\ \hline
Maksym Tomkevych & výskumný pracovník & 00397687 & 2\,000 \\ \hline
\multicolumn{3}{|l|}{\textbf{Spolu: 4 riešitelia}} & \textbf{8\,400} \\ \hline
\multicolumn{3}{|l|}{Ostatní zamestnanci (technická podpora, správa infraštruktúry): 1 osoba} & 600 \\ \hline
\multicolumn{3}{|l|}{\textbf{Celkový počet zamestnancov: 5 \quad Súhrnná kapacita v~hodinách}} & \textbf{9\,000} \\ \hline
\end{tabularx}

% ---------------------------------------------------------------
\begin{formular}{VV-A4}{Základné informácie o zodpovednom riešiteľovi}
\polozka{01}{Meno a priezvisko}{Andrii Tsybulskyi}
\polozka{02}{Organizácia}{STU v~Bratislave, FIIT}
\polozka{03}{Zodpovedný riešiteľ je mladým vedeckým pracovníkom}{áno (do 35 rokov)}
\polozka{04}{Projekty za posledných 5 rokov}{0 -- zodpovedný riešiteľ je na začiatku výskumnej kariéry; projekt je prvým samostatným projektom tímu.}
\polozka{05}{Expertízy, konzultácie a výsledky s~priamym využitím v~praxi}{0}
\polozka{06}{Aplikačné výstupy -- chránené}{0}
\end{formular}

\noindent\textbf{Čestné vyhlásenie zodpovedného riešiteľa.} Ja, dolu podpísaný Andrii Tsybulskyi, čestne vyhlasujem, že údaje uvedené v~častiach ,,Základné informácie o~zodpovednom riešiteľovi`` sú pravdivé a že pri riešení projektu budú dodržané všetky zásady etických princípov vyplývajúce z~európskej a národnej legislatívy.

\medskip\noindent Miesto: Bratislava \hfill Dátum: 1.~11.~2026 \hfill Podpis: \rule{4cm}{0.4pt}

% ---------------------------------------------------------------
\begin{formular}{VV-B}{Ciele, zámery a výstupy projektu}
\polozka{01}{Kľúčové slová}{veľké jazykové modely; AI agenti; penetračné testovanie; bezpečnosť webových aplikácií; benchmark; hodnotenie AI}
\polozka{02}{Ciele projektu}{%
C1: navrhnúť a implementovať Web Pentest Harness -- izolované prostredie na spúšťanie AI agentov proti zraniteľným webovým aplikáciám so záznamom všetkých akcií; \newline
C2: vytvoriť benchmark aspoň 40 úloh v~minimálne 8 triedach zraniteľností podľa OWASP Top 10 s~anotovanou referenčnou databázou; \newline
C3: navrhnúť hodnotiaci modul s~metrikami presnosť, úplnosť, $F_1$, miera falošných nálezov, čas do nálezu a náklady behu; \newline
C4: porovnať aspoň 4 architektúry AI agentov a 3 jazykové modely s~klasickým DAST nástrojom a publikovať výsledky.}
\polozka{03}{Zaradenie projektu do TRL}{TRL 3 $\rightarrow$ TRL 5}
\polozka{03}{Zdôvodnenie zaradenia do TRL}{Princíp autonómneho testovania pomocou LLM agentov je experimentálne overený v~literatúre (TRL 3). Projekt ho overí v~relevantnom, no kontrolovanom prostredí s~realistickými aplikáciami (TRL 4--5).}
\polozka{04}{Zaradenie výstupov do TRL}{Web Pentest Harness (softvér, TRL 5); benchmark a referenčná databáza (dátová sada, TRL 5); metodika hodnotenia (TRL 4).}
\polozka{04}{Využitie výsledkov -- odberateľ je žiadateľ}{áno -- FIIT STU využije harness vo výučbe bezpečnosti a vo výskume}
\polozka{04}{Využitie výsledkov -- iný odberateľ}{áno -- bezpečnostné firmy a CSIRT tímy na porovnanie nástrojov}
\polozka{05}{Prehľad plánovaných verejne dostupných výstupov}{%
E1: technická správa s~návrhom architektúry a taxonómiou úloh; \newline
E2: otvorený zdrojový kód harness a prvá verzia benchmarku; \newline
E3: dátová sada výsledkov experimentov, článok na konferencii; \newline
E4: finálna verzia benchmarku, článok v~časopise, metodická príručka.}
\end{formular}

% ---------------------------------------------------------------
\section*{VV-C \quad {\normalsize Rozpočet projektu v~EUR}}
\renewcommand{\arraystretch}{1.25}
\begin{tabularx}{\textwidth}{|l|X|r|r|r|r|r|}
\hline
\rowcolor{formbg} & \textbf{Položka} & \textbf{2027} & \textbf{2028} & \textbf{2029} & \textbf{2030} & \textbf{Spolu} \\ \hline
01 & Bežné priame náklady & 34\,700 & 69\,400 & 69\,400 & 34\,700 & 208\,200 \\ \hline
02 & \quad Mzdové a ostatné osobné náklady & 20\,000 & 40\,000 & 40\,000 & 20\,000 & 120\,000 \\ \hline
03 & \quad Poistné & 7\,200 & 14\,400 & 14\,400 & 7\,200 & 43\,200 \\ \hline
04 & \quad Cestovné náklady & 1\,500 & 3\,000 & 3\,000 & 1\,500 & 9\,000 \\ \hline
05 & \quad Materiál & 1\,000 & 2\,000 & 2\,000 & 1\,000 & 6\,000 \\ \hline
06 & \quad Služby (cloud, prístup k~API modelov) & 4\,000 & 8\,000 & 8\,000 & 4\,000 & 24\,000 \\ \hline
07 & \quad Energie, komunikácie & 500 & 1\,000 & 1\,000 & 500 & 3\,000 \\ \hline
08 & \quad Náklady na popularizáciu & 500 & 1\,000 & 1\,000 & 500 & 3\,000 \\ \hline
09 & Bežné nepriame náklady & 5\,300 & 10\,600 & 10\,600 & 5\,300 & 31\,800 \\ \hline
10 & \textbf{Bežné náklady spolu = celkové náklady z~APVV} & \textbf{40\,000} & \textbf{80\,000} & \textbf{80\,000} & \textbf{40\,000} & \textbf{240\,000} \\ \hline
 & Spolufinancovanie & 0 & 0 & 0 & 0 & 0 \\ \hline
 & \textbf{Celkové náklady} & \textbf{40\,000} & \textbf{80\,000} & \textbf{80\,000} & \textbf{40\,000} & \textbf{240\,000} \\ \hline
\end{tabularx}

\medskip\noindent Kapitálové náklady sa neplánujú; výpočtový výkon sa zabezpečí ako služba (položka 06).

% ---------------------------------------------------------------
\section*{VV-D \quad {\normalsize Harmonogram projektu}}
\begin{tabularx}{\textwidth}{|c|X|l|l|}
\hline
\rowcolor{formbg}\textbf{Etapa} & \textbf{Názov etapy} & \textbf{Začiatok} & \textbf{Koniec} \\ \hline
E1 & Analýza stavu poznania, návrh architektúry harness a taxonómie úloh & 1.~7.~2027 & 31.~12.~2027 \\ \hline
E2 & Implementácia harness, zraniteľných aplikácií a referenčnej databázy & 1.~1.~2028 & 31.~12.~2028 \\ \hline
E3 & Experimenty s~AI agentmi a klasickými nástrojmi, priebežné vyhodnotenie & 1.~1.~2029 & 31.~12.~2029 \\ \hline
E4 & Finálne vyhodnotenie, otvorené vydanie benchmarku, diseminácia & 1.~1.~2030 & 30.~6.~2030 \\ \hline
\end{tabularx}

\subsection*{Zoznam prístrojov v~projekte}
Projekt neplánuje nákup prístrojov.

% ---------------------------------------------------------------
\section*{VV-E \quad {\normalsize Čestné vyhlásenie štatutárneho zástupcu žiadateľskej organizácie}}
Ja, dolu podpísaný/á, štatutárny zástupca / štatutárna zástupkyňa organizácie Slovenská technická univerzita v~Bratislave, záväzne vyhlasujem, že údaje uvedené v~tejto žiadosti sú pravdivé, že žiadateľ spĺňa podmienky výzvy a že svojím podpisom zodpovedám za správnosť a pravdivosť údajov týkajúcich sa organizácie, za ktorú sa predkladá toto čestné vyhlásenie.

\medskip\noindent Miesto: Bratislava \hfill Dátum: \rule{3cm}{0.4pt} \hfill Podpis štatutárneho zástupcu: \rule{4cm}{0.4pt}

\end{document}
